% !PRINTED 491
\weberpage{491}

But these two permutations do not form a group, because the repetition of the latter,
\[
\begin{pmatrix} 0, & 1, & 2 \\ 2, & 0, & 1 \end{pmatrix}
\]
does not occur among them.

For these functions the following theorem now holds:

5. \gesperrt{If $P$ is a permutation group of $m$ figures, and if $\alpha, \alpha_1 \dots \alpha_{m-1}$ are arbitrary quantities different from one another, then there are rational functions of the $\alpha$, even with rational coefficients, which do not change when a permutation of the group $P$ is applied to the indices of $\alpha$, and which do change when a permutation not belonging to $P$ is applied.}

To prove it, we take a function
\[
\varrho = \Phi(\alpha, \alpha_1 \dots \alpha_{m-1}),
\]
such as we have already considered in \S.\,143, which takes $\Pi(m)$ different values when the $\alpha$ are permuted in all possible ways.

If we denote the values which $\varrho$ takes through the application of the permutations of a group $P$ by
\weq{1}{\varrho, \varrho_1 \dots \varrho_{\mu-1},}
then by Theorem 2. these quantities are only interchanged among one another when a permutation $\pi_1$ from the group $P$ is applied to all of them at once. For the application of the permutation $\pi_1$ to all the functions (1) has the same effect as the application of the permutations contained in $P_1$ to $\varrho$. But by 2., $P_1$ is identical with $P$, apart from the order.

If, then, $t$ denotes a variable, the function
\weq{2}{\psi = (t - \varrho)(t - \varrho_1) \dots (t - \varrho_{\mu-1})}
remains unchanged when any permutation from $P$ is applied. Now if, by a permutation not contained in $P$, the quantities (1) pass into
\weq{3}{\varrho', \varrho_1' \dots \varrho_{\mu-1}'}
then at least $\varrho'$ certainly does not occur among the quantities (1), and consequently
\[
\psi_1 = (t - \varrho')(t - \varrho_1') \dots (t - \varrho_{\mu-1}')
\]
%CONT>>
